
\documentclass{article}
\usepackage{colortbl}
\usepackage{makecell}
\usepackage{multirow}
\usepackage{supertabular}

\begin{document}

\newcounter{utterance}

\centering \large Interaction Transcript for game `cladder', experiment `full\_v1.5\_default', episode 9771 with reflexion{-}llama{-}3.1{-}8b{-}instruct.
\vspace{24pt}

{ \footnotesize  \setcounter{utterance}{1}
\setlength{\tabcolsep}{0pt}
\begin{supertabular}{c@{$\;$}|p{.15\linewidth}@{}p{.15\linewidth}p{.15\linewidth}p{.15\linewidth}p{.15\linewidth}p{.15\linewidth}}
    \# & \multicolumn{2}{c}{Player} && \multicolumn{2}{c}{Game Master} \\
    \hline

    \theutterance \stepcounter{utterance}  
    & & & \multicolumn{4}{p{0.6\linewidth}}{
        \cellcolor[rgb]{0.9,0.9,0.9}{
            \makecell[{{p{\linewidth}}}]{
                \texttt{\tiny{[P1$\langle$GM]}}
                \texttt{You are an expert in causal inference. The following question is not a typical commonsense query, but rather a meticulously designed question created by a professor specializing in causal inference, intended to assess the students' mastery of the course content.} \\
\texttt{Imagine a self{-}contained, hypothetical world with only the following conditions, and without any unmentioned factors or causal relationships: Jyka has a direct effect on yupt and kwox. Yupt has a direct effect on kwox. We know that jyka causes yupt. jyka or yupt causes kwox. Would an individual is kwox if not jyka instead of jyka?} \\
\texttt{Start your answer with 'Yes' or 'No', followed by additional reasoning or evidence to support your explanation.} \\
            }
        }
    }
     \\ \\

    \theutterance \stepcounter{utterance}  
    & \multicolumn{4}{p{0.6\linewidth}}{
        \cellcolor[rgb]{0.561,0.737,0.561}{
            \makecell[{{p{\linewidth}}}]{
                \texttt{\tiny{[P1$\rangle$GM]}}
                \texttt{No.} \\
\\ 
\texttt{To answer this question, we need to apply the principles of causal inference, specifically the concept of backdoor paths and the do{-}calculus. } \\
\\ 
\texttt{First, let's establish the causal graph:} \\
\\ 
\texttt{{-} Jyka → Yupt} \\
\texttt{{-} Jyka → Kwox} \\
\texttt{{-} Yupt → Kwox} \\
\texttt{{-} Jyka or Yupt → Kwox (this is a more complex relationship, but we can interpret it as Jyka → Kwox and Yupt → Kwox, with the possibility of confounding between Jyka and Yupt)} \\
\\ 
\texttt{We are asked to evaluate the effect of not Jyka (i.e., Jyka=0) on Kwox. To do this, we need to consider the backdoor paths from Jyka to Kwox.} \\
\\ 
\texttt{There are two backdoor paths:} \\
\\ 
\texttt{1. Jyka → Yupt → Kwox} \\
\texttt{2. Jyka → Kwox (direct effect)} \\
\\ 
\texttt{To block these backdoor paths, we need to condition on the variables that are on the path. In this case, we need to condition on Yupt to block the first backdoor path.} \\
\\ 
\texttt{Now, let's consider the effect of Jyka=0 on Kwox, conditioning on Yupt. If Jyka=0, then Yupt will be 0 (since Jyka causes Yupt). However, Yupt still has a direct effect on Kwox, so Kwox will be 0 even if Jyka=0.} \\
\\ 
\texttt{Therefore, the answer is No: an individual would still be Kwox if not Jyka instead of Jyka, conditioning on Yupt.} \\
            }
        }
    }
    & & \\ \\

    \theutterance \stepcounter{utterance}  
    & & & \multicolumn{2}{p{0.3\linewidth}}{
        \cellcolor[rgb]{0.9,0.9,0.9}{
            \makecell[{{p{\linewidth}}}]{
                \texttt{\tiny{[GM$|$GM]}}
                \texttt{No} \\
            }
        }
    }
    & & \\ \\

    \theutterance \stepcounter{utterance}  
    & & & \multicolumn{2}{p{0.3\linewidth}}{
        \cellcolor[rgb]{0.9,0.9,0.9}{
            \makecell[{{p{\linewidth}}}]{
                \texttt{\tiny{[GM$|$GM]}}
                \texttt{no} \\
            }
        }
    }
    & & \\ \\

    \theutterance \stepcounter{utterance}  
    & & & \multicolumn{2}{p{0.3\linewidth}}{
        \cellcolor[rgb]{0.9,0.9,0.9}{
            \makecell[{{p{\linewidth}}}]{
                \texttt{\tiny{[GM$|$GM]}}
                \texttt{game\_result = WIN} \\
            }
        }
    }
    & & \\ \\

\end{supertabular}
}

\end{document}
